\documentclass[10pt,a4paper]{article}

\usepackage[T1]{fontenc}
\usepackage[utf8]{inputenc}
\usepackage{textcomp}\usepackage{lmodern,amsmath,amssymb,array,tabularx,booktabs,xcolor,enumitem}
\usepackage[margin=20mm]{geometry}
\usepackage[unicode,colorlinks=true,urlcolor=black,linkcolor=black]{hyperref}
\usepackage{xcolor}
\usepackage{xltabular}
\usepackage[normalem]{ulem}
\usepackage{kotex}
\usepackage{mdframed}

\setlength{\parindent}{0pt}
\setlength{\parskip}{5pt}
\setlength{\emergencystretch}{2em}
\setlength{\tabcolsep}{4pt}
\renewcommand{\arraystretch}{1.10}
\newcommand{\hd}[1]{\par\vspace{8pt}\textbf{#1}\par\vspace{3pt}}
\newcommand{\circnum}[1]{\textcircled{\scriptsize #1}}
\newcommand{\note}{\textreferencemark\ }
\newcommand{\red}[1]{\textcolor{red}{#1}}
\newcommand{\under}[2]{\underbrace{#1}_{\text{#2}}}
\newcommand{\paper}[1]{\href{#1}{paper}}
\newcolumntype{Y}{>{\raggedright\arraybackslash}X}

\begin{document}


\section{Driving factors}
This document shows the summary of candidate driving factors of \textbf{the number of potential ride-hailing usage opportunities($M$)} and \textbf{probability of choosing robotaxi for that
usage opportunity} derived from the documents of \texttt{02\_demand\_variables} and \texttt{03\_city\_entry\_conditions}.

\hd{A. Driving factors of $M$}
The papers below mainly explain \textbf{observed ride-hailing usage}, so they do not include the `unrealized demand' included in the definition of $M$, making direct estimation impossible. However, I suppose the demand patterns provided by these papers lead to constructing `potential travel opportunities $M$' through modeling.

{\small
\begin{tabularx}{\linewidth}{@{}>{\raggedright\arraybackslash}p{0.16\linewidth}>{\raggedright\arraybackslash}p{0.32\linewidth}>{\raggedright\arraybackslash}p{0.3\linewidth}Y@{}}
\toprule
Variable & Supporting paper & Scope established in the paper & Measurement data \\ \midrule
\textbf{Population and population density} & Elkhouly and Alhadidi, \textit{Ride-sharing Determinants} (\paper{https://arxiv.org/abs/2406.11590}); Zhang et al., \textit{Examining spatial heterogeneity\ldots} (\paper{https://arxiv.org/abs/2209.07980}) & Relationship of population and local density with conventional ride-hailing usage. Density effects may vary by area & \href{https://www.census.gov/programs-surveys/acs}{ACS} \\
\textbf{Employment and economic activity} & Elkhouly and Alhadidi (\paper{https://arxiv.org/abs/2406.11590}) & Uses the economically active population and employment environment to explain usage & \href{https://lehd.ces.census.gov/data/}{LEHD/LODES}, ACS \\
\textbf{Income} & Elkhouly and Alhadidi (\paper{https://arxiv.org/abs/2406.11590}) & Analyzes the relationship between income and usage in Chicago & \href{https://data.census.gov/table/ACSDT5Y2024.B19013}{ACS B19013} \\
\textbf{Share of car-free households} & Elkhouly and Alhadidi (\paper{https://arxiv.org/abs/2406.11590}) & Analyzes the relationship between households without cars and usage & \href{https://data.census.gov/table/ACSDT1Y2022.B08201}{ACS B08201} \\
\textbf{Land-use mix, accessibility, and commercial activity} & Zhang et al., \textit{Examining spatial heterogeneity\ldots} (\paper{https://arxiv.org/abs/2209.07980}) & Nonlinear and spatial differences in the urban environment and ridesourcing usage & \href{https://www.epa.gov/smartgrowth/smart-location-mapping}{EPA SLD} \\
\textbf{Restaurant and leisure-facility density} & Loa, Lee, and Circella, \textit{Just another ride-hailing?} (\paper{https://link.springer.com/article/10.1007/s11116-026-10751-9}) & Relationship with \textbf{robotaxi trip counts in SF}. Does not separate the effects of $M$ and $p$ & \href{https://www.census.gov/programs-surveys/cbp.html}{Census CBP}, OpenStreetMap \\

\textbf{Public transport and walking environment} & Zhang et al. (\paper{https://arxiv.org/abs/2209.07980}); Cats et al., \textit{Beyond the dichotomy\ldots} (\paper{https://doi.org/10.1371/journal.pone.0262496}) & Analyzes the local environment and travel-time competitiveness and substitutability with public transport. Not a uniform increase/decrease effect & EPA SLD, \href{https://gtfs.org/documentation/overview/}{GTFS} \\
\textbf{Airport-related market size} & Zhang et al. (\paper{https://arxiv.org/abs/2209.07980}) & Analyzes airport travel as a separate spatial category. \textbf{Does not estimate the direct effect of airport passenger counts} & \href{https://www.faa.gov/airports/planning_capacity/passenger_allcargo_stats/passenger}{FAA passenger data}. Used as a proxy for market size \\
\textbf{Time, day of week, and season} & Elkhouly and Alhadidi (\paper{https://arxiv.org/abs/2406.11590}) & Analyzes differences in usage over time. Specific variables such as holidays are selected through separate design & Dates and times in trip data \\
\textbf{Weather} & Additional: Liu, Jiang, and Chen, \textit{Quantifying the impact of weather on ride-hailing ridership: Evidence from Haikou, China} (2021, \paper{https://doi.org/10.1016/j.tbs.2021.04.002}) & Analyzes the relationship between weather and ride-hailing usage. Effects in U.S. cities require revalidation & \href{https://www.ncei.noaa.gov/cdo-web/}{NOAA CDO} \\
\bottomrule
\end{tabularx}
}

%The determinants of $M$ are city characteristics (original proposal) + temporal/seasonal context (situation) + market activity (outcome: past Uber/Lyft usage) Past usage is a predictor that can be used separately. However, this does not explain ``why demand occurs''; it can be classified as \textbf{outcome data showing the size (scale) of demand.}

%------------------------------
\newpage
\textit{\textbf{\textcolor{blue}
{[Revised on 10/01/2026]}}}
\begin{mdframed}[linecolor=red!60,linewidth=2pt]
\color{gray!60}


\hd{B. Driving factors of $p$}

{\small
\begin{tabularx}{\linewidth}{@{}>{\raggedright\arraybackslash}p{0.16\linewidth}>{\raggedright\arraybackslash}p{0.32\linewidth}>{\raggedright\arraybackslash}p{0.3\linewidth}Y@{}}
\toprule
Variable & Supporting paper & Scope of application & Required measurement data \\ \midrule
\textbf{Relative fare} & Additional: Yin and Cherchi, \textit{Willingness to pay for automated taxis\ldots} (2024; online 2022, \paper{https://doi.org/10.1007/s11116-022-10319-3}) & Choice experiment between conventional taxis and automated taxis & Offered fares for both + choice outcomes \\
\textbf{Relative waiting time} & Yin and Cherchi (\paper{https://doi.org/10.1007/s11116-022-10319-3}) & Analyzes choices under different service conditions & Offered waiting times for both + choice outcomes \\
\textbf{In-vehicle travel time} & Yin and Cherchi (\paper{https://doi.org/10.1007/s11116-022-10319-3}) & Analyzes choice trade-offs between fares and time & Offered travel times for both \\
\textbf{Trust, price value, and usage habits} & Additional: Wei et al., \textit{Using the Extended Unified Theory\ldots} (2024, \paper{https://www.nature.com/articles/s41599-024-03271-3}) & Study of robotaxi acceptance intention and behavior. \textbf{Survey scales cannot be interpreted directly as choice probabilities} & Individual attitude and habit surveys \\
\textbf{Age and robotaxi-related knowledge} & Yin and Cherchi (\paper{https://doi.org/10.1007/s11116-022-10319-3}) & Choice differences by individual characteristics. Fixing the sign, such as ``younger people always prefer it, is inappropriate & Link individual characteristics to choice data \\
\textbf{User reviews and in-vehicle support features} & Yin and Cherchi (\paper{https://doi.org/10.1007/s11116-022-10319-3}) & Analyzes customer ratings and features such as destination changes and communication with an operator as choice attributes & Service attributes in the choice experiment \\
\textbf{Experience using robotaxis} & Additional: Yin, Cherchi, and Nettle, \textit{Living in a Ci ty Where Automated Taxis are Operating and Using Them\ldots} (2023, \paper{https://doi.org/10.1177/03611981231155181}) & Analyzes preference differences by whether service operates in the residential area and by individual usage experience & Individual usage experience \\

\textbf{Competitiveness of alternatives such as public transport} & Cats et al., \textit{Beyond the dichotomy\ldots} (\paper{https://doi.org/10.1371/journal.pone.0262496}) & Evidence from \textbf{conventional ride-hailing}. Not a robotaxi-specific preference effect & Availability, fares, and times of alternatives \\
\bottomrule
\end{tabularx}
}

\end{mdframed}


%------------revised version-----
\newpage
\textbf{\textcolor{blue}{[Revised version]}}
%%%%%%%%%%%%%%%%%%%%%%%%%%%%%%%%%%%%%%%%
\hd{B. Driving factors of $p$}

{\small
\begin{xltabular}{\linewidth}
{@{}>{\raggedright\arraybackslash}p{0.16\linewidth}
>{\raggedright\arraybackslash}p{0.32\linewidth}
>{\raggedright\arraybackslash}p{0.3\linewidth}Y@{}}
\toprule
Variable & Supporting paper & Scope of application
& Required measurement data \\
\midrule

\textbf{Relative fare}
& Yin and Cherchi,
\textit{Willingness to pay for automated taxis}
(\paper{https://doi.org/10.1007/s11116-022-10319-3})
& Choice experiment between conventional and automated taxis.
Examines responses to fare differences.
Own-service fare depends on operator decisions.
& Fares of competing services + choice outcomes \\

\textbf{Relative waiting time}
& Yin and Cherchi
(\paper{https://doi.org/10.1007/s11116-022-10319-3})
& Examines responses to waiting-time differences.
Own-service waiting time depends on fleet size and allocation.
& Waiting times of competing services + choice outcomes \\

\textbf{In-vehicle travel time}
& Yin and Cherchi
(\paper{https://doi.org/10.1007/s11116-022-10319-3})
& Examines trade-offs between fare and travel time.
Distinguish time affected by operating policies
from underlying traffic conditions.
& Travel times of competing services + choice outcomes \\

\textbf{Trust, price value, and usage habits}
& Wei et al.,
\textit{Using the Extended Unified Theory}
(\paper{https://www.nature.com/articles/s41599-024-03271-3});
Yin, Cherchi, and Nettle (2023)
& Evidence on acceptance intentions and behavior related
to price value, habits, and experience.
Wei et al. do not include trust.
Survey scores are not choice probabilities.
& Individual attitude and habit surveys;
choice and usage data \\

\textbf{Age and robotaxi-related knowledge}
& Yin and Cherchi
(\paper{https://doi.org/10.1007/s11116-022-10319-3});
Wei et al.
(\paper{https://www.nature.com/articles/s41599-024-03271-3})
& Choice and acceptance differences by individual characteristics.
Do not impose a uniform sign on age.
Distinguish general AV knowledge from awareness of actual service.
& Individual age, knowledge, and service awareness
linked to choice data \\

\textbf{User reviews and in-vehicle support features}
& Yin and Cherchi
(\paper{https://doi.org/10.1007/s11116-022-10319-3})
& Examines reviews, destination changes, and operator contact.
Support features are potential operator decisions;
reviews also reflect past usage.
& Reviews and service features + choice outcomes \\

\textbf{Experience using robotaxis}
& Yin, Cherchi, and Nettle,
\textit{Living in a City Where Automated Taxis
are Operating and Using Them}
(\paper{https://doi.org/10.1177/03611981231155181})
& Relates preferences to residence in an operating area
and individual usage experience.
A uniform positive experience effect is not established.
& Individual experience before the decision;
service-operation records \\

\textbf{Competitiveness of alternatives}
& Cats et al.,
\textit{Beyond the dichotomy}
(\paper{https://doi.org/10.1371/journal.pone.0262496})
& Evidence from conventional ride-hailing.
Substitution and complementarity depend on
travel-time competitiveness with alternatives.
& Alternative availability, fares, and travel times;
GTFS, NTD \\

%---B-2. Urban and socioeconomic characteristics}-------------


\textbf{Share of car-free households}
& Elkhouly and Alhadidi,
\textit{Ride-sharing Determinants}
(\paper{https://arxiv.org/abs/2406.11590});
Loa, Lee, and Circella,
\textit{Just another ride-hailing?}
(\paper{https://doi.org/10.1007/s11116-026-10751-9})
& Positive associations with conventional ride-hailing
usage in Chicago and robotaxi trip origins in SF.
Candidate measure of limited access to a private-car alternative.
& \href{https://www.census.gov/programs-surveys/acs}{ACS}
B08201: car-free households / all households \\

\textbf{Income}
& Elkhouly and Alhadidi;
Loa, Lee, and Circella;
Yin and Cherchi;
Wei et al.
& Positive associations in aggregate usage data
from Chicago and SF.
Individual choice experiments and surveys
do not consistently establish the same effect.
& ACS B19013; income-group shares if needed \\

\textbf{Employment and economic activity}
& Elkhouly and Alhadidi
(\paper{https://arxiv.org/abs/2406.11590});
Loa, Lee, and Circella
(\paper{https://doi.org/10.1007/s11116-026-10751-9})
& Economic activity share is positively associated
with conventional ride-hailing usage in Chicago;
employment density with robotaxi trip destinations in SF.
Distinguish the two measures.
& \href{https://lehd.ces.census.gov/data/}{LEHD/LODES};
ACS B23025 \\

\textbf{Population density}
& Elkhouly and Alhadidi;
Zhang et al.,
\textit{Examining spatial heterogeneity}
(\paper{https://arxiv.org/abs/2209.07980});
Loa, Lee, and Circella
& Spatially varying relationships for conventional ride-hailing.
Negative association with robotaxi usage in SF.
Effects at the city level require separate validation.
& ACS population + area within the same boundary \\

\textbf{Land-use mix, accessibility, and commercial activity}
& Zhang et al.
(\paper{https://arxiv.org/abs/2209.07980});
Loa, Lee, and Circella
(\paper{https://doi.org/10.1007/s11116-026-10751-9})
& Nonlinear and spatially heterogeneous relationships
for conventional ride-hailing.
Positive association with land-use mix in SF.
Select a representative measure and aggregation method.
& \href{https://www.epa.gov/smartgrowth/smart-location-mapping}
{EPA Smart Location Database} \\

\textbf{Public transport and walking environment}
& Zhang et al.;
Cats et al.
(\paper{https://doi.org/10.1371/journal.pone.0262496});
Loa, Lee, and Circella
& Substitution and complementarity coexist
for conventional ride-hailing.
Positive associations of transit frequency and walkability
with robotaxi usage in SF.
& EPA SLD; GTFS; NTD \\

\textbf{Restaurant and leisure-facility density}
& Loa, Lee, and Circella,
\textit{Just another ride-hailing?}
(\paper{https://doi.org/10.1007/s11116-026-10751-9})
& Positive association with robotaxi trip counts in SF.
The evidence does not separate the effects of $M$ and $p$.
& \href{https://www.census.gov/programs-surveys/cbp.html}
{Census CBP}; OpenStreetMap \\

%------- Additional factors---

\textbf{Age composition}
& Loa, Lee, and Circella;
Etminani-Ghasrodashti et al. (2022)
(\paper{https://doi.org/10.1177/03611981221093632})
& Aggregate and pilot evidence concerning younger users.
Individual age effects do not directly establish
effects of an area's age composition.
& ACS B01001; initially consider the share aged 18--34 \\

\textbf{Student share}
& Etminani-Ghasrodashti et al. (2022)
(\paper{https://doi.org/10.1177/03611981221093632})
& Actual usage evidence from the Arlington campus pilot.
Account for free student rides and the campus setting.
& ACS enrollment data; local student population;
fare-support conditions \\

\textbf{Education and technology affinity}
& Krueger, Rashidi, and Rose (2016);
Haboucha et al. (2017);
Gkartzonikas and Gkritza (2019)
& Additional candidates considered in SAV
choice experiments and reviews.
Education and technology affinity are distinct;
specific findings require confirmation.
& ACS education data; technology-affinity surveys \\

\textbf{Existing ride-hailing habits}
& Wei et al.
(\paper{https://www.nature.com/articles/s41599-024-03271-3});
conventional ride-hailing studies
& Examines relationships of habits and previous usage
with adoption.
Past trip counts are a proxy for habits,
not the same measure as a habit scale.
& Common habit surveys;
conventional ride-hailing usage before the decision \\

\textbf{Safety perceptions and exposure}
& SAV choice experiments and reviews;
bibliographic details and findings of
Khan et al. (2022) require confirmation
& Distinguish perceived safety from actual crash or crime exposure.
Post-use satisfaction should not directly substitute
for safety perceptions before the decision.
& Common safety surveys;
crash, crime, and exposure data \\

\textbf{Regulatory permission}
& Permits and operating requirements
from the relevant authorities;
NCSL legislative information
& An operating-feasibility condition.
Treat areas without the required permission
separately in entry feasibility.
& Valid permits and operating requirements by area \\

\textbf{Climate and ODD}
& Weather records and operators'
Operational Design Domain information
& Climate is an exogenous condition.
ODD and service availability depend on technical
and operating conditions; distinguish these from preference effects.
& \href{https://www.ncei.noaa.gov/cdo-web/}{NOAA};
operator ODD; operating-restriction records \\

\textbf{Time since service launch}
& Yin, Cherchi, and Nettle (2023)
(\paper{https://doi.org/10.1177/03611981231155181});
service-launch and operation records
& Candidate proxy for accumulated awareness and experience.
Launch timing depends on operator decisions;
the effect of elapsed time requires separate validation.
& Company announcements; permit and launch dates;
annual operation records \\

\bottomrule
\end{xltabular}
}

\textbf{Measurement principle.}
Robotaxi trips and $M$ must use the same geographic boundary,
year, and trip unit.
If the numerator contains completed trips,
$p$ measures a realized usage share affected by supply conditions.

\end{document}